
\subsubsection{6.1 Key Findings}

\textbf{Contracts and differential oracles measure fundamentally different properties.} The 40\% contract-unique mass demonstrates that ''more tests'' cannot close the oracle-strength gap: programs that output the correct answer on CVT inputs --- thereby satisfying the differential oracle --- nonetheless violate reference contracts on those same inputs in 40\% of cases. This gap is not marginal; it is $4\times$ the predicted threshold and holds across all 364 tasks and both task types.

\textbf{The oracle-isolation gap reflects a threshold effect of contract presence.} Spearman $\rho$ = 0.136 and the non-monotonic tier gradient (T2 > T4 > T3 > T1) indicate that the presence of any formal pre/post-condition assertion --- even a simple input validity check --- creates a substantial oracle-isolation gap (~0.40) on CVT inputs. The additional contribution of greater contract complexity is weak and does not follow a monotone gradient. This finding likely reflects ContractEval's precondition-dominated contract structure: Tier 2 BoolOp contracts function as strict input guards, making them highly likely to fire on CVT inputs; the AST-tier system does not cleanly separate precondition strictness from postcondition semantic richness.

\textbf{Adaptive PBT provides a model-independent evaluation enhancement.} The 0.099--0.103 contribution range across 5 diverse model families argues for adaptive PBT as a property of the contract oracle and input exploration space, not a model-specific characteristic. The icontract-hypothesis adaptive contribution is consistent across open-source and closed-source models, large and small parameter counts, and both task types.

\textbf{Cross-model contract-satisfaction variation is negligible within this evaluation scope.} The $\Delta R^2$ = 0.004 and task-controlled residual gap of 0.007 indicate that, after controlling for pass@1* and model size, model-family identity adds negligible explanatory power in the 5-model set evaluated. Task difficulty accounts for substantially more between-unit variance than model identity in this scope.

\subsubsection{6.2 Limitations}

\textbf{L1: CVT input scope.} The oracle-isolation gap is measured on contract-violating inputs, not on a general or random input distribution. CVT inputs are the natural evaluation class for contract-violating behavior --- they are the inputs that formal contracts are designed to discriminate. Whether the 40\% oracle gap persists on random input distributions is an open question. All quantitative claims are explicitly scoped to CVT inputs.

\textbf{L2: Cross-model underpowering.} With n = 5 models, the exact permutation test structurally cannot achieve p < 0.05 for any $\tau$ < 1.0 in a two-sided test. The $\Delta R^2$ = 0.004 finding is interpretable as a genuine null independent of statistical power: model family adds negligible variance beyond capability proxies. Confirming or refuting the $\tau$ direction requires n $\geq$ 10 models.

\textbf{L3: ContractEval scope.} All findings are scoped to 364 Python algorithmic tasks with inline assert contracts. Generalization to software engineering tasks, non-Python languages, or other contract formats (icontract decorators, Dafny, SPARK) is an open direction.

\textbf{L4: Richness gradient null.} The $\rho$ = 0.136 finding likely reflects precondition-dominated contract structure; the richness tier system captures syntactic AST complexity rather than oracle discriminating power. Filtering to postcondition-only contracts may recover a stronger gradient, as HumanEval+ tasks --- which have higher $\rho$ (0.164) than MBPP+ (0.096) --- tend to have richer postconditions.

\subsubsection{6.3 Relationship to Prior Work}

The oracle-isolation gap of 0.40 extends the EvalPlus result: where EvalPlus showed that dense differential testing reduces pass@1 by up to 23.1\% over sparse testing, this work shows that even 764-test differential evaluation misses 40\% of contract violations on contract-relevant inputs. The adaptive PBT contribution of ~10pp is directionally consistent with Bose [Bose2025Prompts] (18--32\% additional failures on StarCoder/CodeLlama), though not directly comparable --- the present work measures additional violations beyond static CVT inputs rather than beyond unit tests. The execution-based approach closes the tractability gap in ContractEval [Lim2025ContractEval] (25.82\% tractable via Z3 vs. 100\% tractable via execution).

\subsubsection{6.4 Broader Implications}

This work provides empirical support for oracle-type diversity as a first-class dimension in LLM code evaluation. The structural oracle gap --- a property of contract semantics, not input sampling density --- cannot be closed by adding more differential tests. Evaluation frameworks that report only test-pass metrics may systematically classify semantically incorrect programs as correct. The oracle isolation design introduced here provides a reproducible methodology for quantifying this gap on any benchmark with formal contract annotations.

